\documentclass[reqno,11pt]{amsart}
\usepackage{math327-lecture}
\usepackage{tikz}
\usepackage{needspace}
\tcolorboxenvironment{example}{lecturebox,colback=definitiongreen,colframe=definitionborder}
\definecolor{remarkpurple}{HTML}{F2ECFA}
\definecolor{remarkpurpleborder}{HTML}{A68CC7}
\tcolorboxenvironment{remark}{lecturebox,colback=remarkpurple,colframe=remarkpurpleborder}
\lectureinfo{15}{Oct. 14, 2026}{The Correspondence Theorem and Dihedral Groups}

% Transcribed from handwritten Lecture 11, Sep. 26, 2025, pp. 1--12.
% Historical course announcements are omitted. Products, announced on p. 1,
% are developed in the next handwritten lecture and in typed Lecture 16.
\begin{document}
\raggedbottom
\makelecturetitle

\section{Some motivation}

The next result we will study is quite abstract. It lets us learn about
the structure of a group $G$ using a homomorphism from $G$ to another
group. We begin with an example of a surjective homomorphism from
$S_4$ to $S_3$.

There are exactly three ways to partition $\{1,2,3,4\}$ into two
unordered pairs:
\[
\begin{aligned}
	\Pi_1&=\bigl\{\{1,2\},\{3,4\}\bigr\},\\
	\Pi_2&=\bigl\{\{1,3\},\{2,4\}\bigr\},\\
	\Pi_3&=\bigl\{\{1,4\},\{2,3\}\bigr\}.
\end{aligned}
\]
Every $\sigma\in S_4$ permutes these three partitions by applying
$\sigma$ to each entry. Identifying $\Pi_i$ with $i$ gives a map
\[
\varphi:S_4\longrightarrow S_3.
\]
This is a homomorphism: applying $\tau$ and then $\sigma$ to the
partitions has the same effect as applying $\sigma\tau$.

\begin{remark}[A picture of $\varphi$]
Place $1,2,3,4$ at the vertices of a square. The three partitions are
the three ways to join the vertices in two disjoint pairs. The colored
segments below show those pairs:
\begin{center}
\begin{tikzpicture}[x=1cm,y=1cm,
  every node/.style={font=\small},pair/.style={line width=1.5pt}]
  \foreach \offset/\pairname/\paircolor in
    {0/1/courseblue,3.5/2/ForestGreen,7/3/BrickRed} {
    \begin{scope}[xshift=\offset cm]
      \draw[gray!60,thin] (0,0)--(0,1.7)--(1.7,1.7)--(1.7,0)--cycle;
      \ifnum\pairname=1
        \draw[pair,\paircolor] (0,1.7)--(1.7,1.7) (0,0)--(1.7,0);
      \fi
      \ifnum\pairname=2
        \draw[pair,\paircolor] (0,0)--(1.7,1.7) (0,1.7)--(1.7,0);
      \fi
      \ifnum\pairname=3
        \draw[pair,\paircolor] (0,0)--(0,1.7) (1.7,0)--(1.7,1.7);
      \fi
      \fill (1.7,1.7) circle (1.5pt) node[above right] {$1$};
      \fill (0,1.7) circle (1.5pt) node[above left] {$2$};
      \fill (0,0) circle (1.5pt) node[below left] {$3$};
      \fill (1.7,0) circle (1.5pt) node[below right] {$4$};
      \node[\paircolor] at (0.85,-0.65) {$\Pi_{\pairname}$};
    \end{scope}
  }
  \node at (0.85,-1.1) {$\{\{1,2\},\{3,4\}\}$};
  \node at (4.35,-1.1) {$\{\{1,3\},\{2,4\}\}$};
  \node at (7.85,-1.1) {$\{\{1,4\},\{2,3\}\}$};
\end{tikzpicture}
\end{center}
In $\Pi_2$, the pairs are opposite vertices; $\Pi_1$ and $\Pi_3$ use
opposite edges. Any $\sigma\in S_4$ relabels the vertices: if
$\sigma(\Pi_i)=\Pi_j$, then $\varphi(\sigma)(i)=j$.
For the $90^\circ$ counterclockwise rotation $\sigma=(1\,2\,3\,4)$,
\[
  \textcolor{courseblue}{\Pi_1}\longmapsto\textcolor{BrickRed}{\Pi_3},
  \qquad
  \textcolor{ForestGreen}{\Pi_2}\longmapsto\textcolor{ForestGreen}{\Pi_2},
  \qquad
  \textcolor{BrickRed}{\Pi_3}\longmapsto\textcolor{courseblue}{\Pi_1}.
\]
so $\varphi(\sigma)=(1\,3)$.
\end{remark}

\Needspace{7\baselineskip}
\begin{example}
	For $\sigma=(1\,2\,3\,4)$, we obtain
	\[
	\Pi_1\longmapsto\Pi_3,\qquad
	\Pi_2\longmapsto\Pi_2,\qquad
	\Pi_3\longmapsto\Pi_1.
	\]
	Thus $\varphi((1\,2\,3\,4))=(1\,3)$.
	Similarly, $\varphi((1\,2))=(2\,3)$. Since $(1\,3)$ and $(2\,3)$
	generate $S_3$, the map $\varphi$ is surjective.
\end{example}

\begin{example}
For $\sigma=(1\,2\,3)$, the vertex $4$ is fixed. Applying $\sigma$ to
each entry of the partitions gives
\[
\begin{aligned}
  \Pi_1&=\bigl\{\{1,2\},\{3,4\}\bigr\}
    \longmapsto\bigl\{\{2,3\},\{1,4\}\bigr\}=\Pi_3,\\
  \Pi_2&=\bigl\{\{1,3\},\{2,4\}\bigr\}
    \longmapsto\bigl\{\{2,1\},\{3,4\}\bigr\}=\Pi_1,\\
  \Pi_3&=\bigl\{\{1,4\},\{2,3\}\bigr\}
    \longmapsto\bigl\{\{2,4\},\{3,1\}\bigr\}=\Pi_2.
\end{aligned}
\]
Thus the induced permutation cycles the partitions as
$\Pi_1\mapsto\Pi_3\mapsto\Pi_2\mapsto\Pi_1$, and hence
\[
  \varphi((1\,2\,3))=(1\,3\,2).
\]
\end{example}

\Needspace{18\baselineskip}
Recall the \emph{counting formula for homomorphisms} from Lecture~13:
if $\psi:A\to B$ is a homomorphism and $A$ is finite, then
\[
  |A|=|\ker(\psi)|\,|\im(\psi)|.
\]
In particular, if $\psi$ is surjective, then
$|A|=|\ker(\psi)|\,|B|$. We now apply this formula to the surjective
homomorphism $\varphi:S_4\to S_3$.

\begin{example}[The kernel]
Each of the permutations
\[
1,\quad (1\,2)(3\,4),\quad (1\,3)(2\,4),\quad (1\,4)(2\,3)
\]
fixes all three partitions.

The counting formula gives
$|\ker(\varphi)|=24/6=4$, so these are precisely the elements of the
kernel:
\[
K=\ker(\varphi)
=\{1,(1\,2)(3\,4),(1\,3)(2\,4),(1\,4)(2\,3)\}.
\]
This is the \emph{Klein four group} inside $S_4$.
Each of its three nonidentity elements is a product of two disjoint
transpositions, so its square is the identity and its order is $2$.
The identity has order $1$. A cyclic group of order $4$ must have a
generator of order $4$, but $K$ has no such element. Therefore $K$
is not cyclic.
\end{example}

\section{The correspondence theorem}

Let $\varphi:G\to\mathcal G$ be a \emph{surjective} homomorphism, and set
$K=\ker(\varphi)\trianglelefteq G$.
For any subset $\mathcal H\subseteq\mathcal G$, its \emph{preimage}
under $\varphi$ is
\[
  \varphi^{-1}(\mathcal H)=\{g\in G:\varphi(g)\in\mathcal H\}.
\]
This notation does not require $\varphi$ to have an inverse function.

\begin{theorem}[Correspondence theorem]\label{thm:correspondence}
\leavevmode
\begin{enumerate}[label=\textbf{(\arabic*)}]
  \item There is a bijective correspondence
  \[
    \{\text{subgroups } H\leq G \text{ which contain } K \}\quad\longleftrightarrow\quad
    \{\text{subgroups } \mathcal H\leq\mathcal G\},
  \]
  given by the mutually inverse maps
  \[
    H\longmapsto\varphi(H),\qquad
    \mathcal H\longmapsto\varphi^{-1}(\mathcal H).
  \]
  \item Under this correspondence,
  $H\trianglelefteq G$ if and only if
  $\varphi(H)\trianglelefteq\mathcal G$.
  \item If $H$ is finite, then $|H|=|K|\,|\varphi(H)|$.
\end{enumerate}
\end{theorem}


\subsection{The maps are well defined}

For $H\leq G$ containing $K$, the image $\varphi(H)$ is a subgroup of
$\mathcal G$: it is the image of the homomorphism
$\varphi|_H:H\to\mathcal G$.

Conversely, let $\mathcal H\leq\mathcal G$. Then
$\varphi^{-1}(\mathcal H)$ is a subgroup of $G$. Indeed,
$\varphi(e_G)=e_{\mathcal G}\in\mathcal H$, and if
$a,b\in\varphi^{-1}(\mathcal H)$, then
\[
  \varphi(ab)=\varphi(a)\varphi(b)\in\mathcal H.
\]
Also $K\leq\varphi^{-1}(\mathcal H)$, since every element of $K$ maps to
$e_{\mathcal G}$.

\begin{exercise}
Complete the subgroup verification by showing that
$a\in\varphi^{-1}(\mathcal H)$ implies
$a^{-1}\in\varphi^{-1}(\mathcal H)$.
\end{exercise}

\subsection{The maps are inverse to each other}

We claim that if $K\leq H\leq G$, then
\begin{equation}\label{eq:preimage-image}
  \varphi^{-1}(\varphi(H))=H.
\end{equation}
The inclusion $H\subseteq\varphi^{-1}(\varphi(H))$ is immediate.
For the reverse inclusion, suppose
$x\in\varphi^{-1}(\varphi(H))$. Then $\varphi(x)=\varphi(a)$ for some
$a\in H$, so
\[
  \varphi(xa^{-1})=e_{\mathcal G}.
\]
Thus $xa^{-1}\in K\leq H$, and $x=(xa^{-1})a\in H$.

\begin{exercise}
Use the surjectivity of $\varphi$ to prove that
$\varphi(\varphi^{-1}(\mathcal H))=\mathcal H$ for every
$\mathcal H\leq\mathcal G$. Together with
\eqref{eq:preimage-image}, this proves part~(1).
\end{exercise}

\subsection{Normality and orders}

Suppose $H\trianglelefteq G$, and put $\mathcal H=\varphi(H)$.
Given $x\in\mathcal G$, surjectivity gives $x=\varphi(g)$ for some
$g\in G$. Hence
\[
  x\mathcal Hx^{-1}
  =\varphi(g)\varphi(H)\varphi(g)^{-1}
  =\varphi(gHg^{-1})
  =\varphi(H)=\mathcal H.
\]
Therefore $\mathcal H\trianglelefteq\mathcal G$.

\begin{exercise}
Prove the converse: if $\mathcal H\trianglelefteq\mathcal G$, then
$H=\varphi^{-1}(\mathcal H)\trianglelefteq G$.
\end{exercise}

Finally, the restriction $\varphi|_H:H\to\mathcal H$ is surjective,
and its kernel is $H\cap K=K$. The counting formula from Lecture~13
therefore gives
\[
  |H|=|\ker(\varphi|_H)|\,|\im(\varphi|_H)|=|K|\,|\mathcal H|,
\]
when $H$ is finite. This proves part~(3).



\section{Application to \texorpdfstring{$\varphi:S_4\to S_3$}{the map from S4 to S3}}

Write $S_3=\{1,a,a^2,b,ab,a^2b\}$, where $a^3=b^2=1$ and
$bab^{-1}=a^{-1}$. Its subgroups are
\[
  \{1\},\quad \langle a\rangle,\quad
  \langle b\rangle,\quad \langle ab\rangle,\quad
  \langle a^2b\rangle,\quad S_3.
\]
Concretely, take $a=(1\,2\,3)$ and $b=(1\,3)$. Then
$a^2=(1\,3\,2)$, $ab=(2\,3)$, and $a^2b=(1\,2)$, with permutation
products composed from right to left. The subgroup inclusion graph is
shown below; inclusions run upward.
\begin{center}
\begin{tikzpicture}[x=1cm,y=1cm,every node/.style={inner sep=4pt,align=center}]
  \node (top) at (0,3.2) {$S_3$};
  \node (cyclic) at (-4.5,1.6) {$\langle a\rangle$\\$\{1,a,a^2\}$};
  \node (b) at (-1.5,1.6) {$\langle b\rangle$\\$\{1,b\}$};
  \node (ab) at (1.5,1.6) {$\langle ab\rangle$\\$\{1,ab\}$};
  \node (aab) at (4.5,1.6) {$\langle a^2b\rangle$\\$\{1,a^2b\}$};
  \node (bottom) at (0,0) {$\{1\}$};
  \draw[courseblue,thick] (bottom) -- (cyclic) -- (top);
  \draw[courseblue,thick] (bottom) -- (b) -- (top);
  \draw[courseblue,thick] (bottom) -- (ab) -- (top);
  \draw[courseblue,thick] (bottom) -- (aab) -- (top);
\end{tikzpicture}
\end{center}

Consequently, there are exactly six subgroups of $S_4$ containing $K$.
Their orders follow from $|\varphi^{-1}(\mathcal H)|=4|\mathcal H|$:
\[
\begin{array}{c|c|c}
  \mathcal H\leq S_3 & \varphi^{-1}(\mathcal H)\leq S_4
    & |\varphi^{-1}(\mathcal H)|\\ \hline
  \{1\} & K & 4\\
  \langle a\rangle=\langle(1\,2\,3)\rangle & A & 12\\
  \langle b\rangle=\langle(1\,3)\rangle & H_1 & 8\\
  \langle ab\rangle=\langle(2\,3)\rangle & H_2 & 8\\
  \langle a^2b\rangle=\langle(1\,2)\rangle & H_3 & 8\\
  S_3 & S_4 & 24
\end{array}
\]

What are the groups $A,H_1,H_2,H_3$?

First, we identify $A$ with $A_4$. Every element of $K$ is an even
permutation, so $K\leq A_4$. The restriction of $\varphi$ to $A_4$
therefore has kernel $K$, and the counting formula gives
\[
  |\varphi(A_4)|=\frac{|A_4|}{|K|}=\frac{12}{4}=3.
\]
The unique subgroup of order $3$ in $S_3$ is $\langle a\rangle$.
Thus $\varphi(A_4)=\langle a\rangle$, and the correspondence theorem
gives
\[
  A=\varphi^{-1}(\langle a\rangle)=A_4.
\]

The three subgroups $H_i$ of order $8$ are examples of
\emph{dihedral groups}, as we now explain.

Label the vertices of a square consecutively as follows:
\begin{center}
\begin{tikzpicture}[scale=1.15]
  \draw[thick] (0,0)--(0,2)--(2,2)--(2,0)--cycle;
  \draw[dashed,courseblue,thick] (-0.35,-0.35)--(2.35,2.35);
  \fill (2,2) circle (1.7pt) node[above right] {$1$};
  \fill (0,2) circle (1.7pt) node[above left] {$2$};
  \fill (0,0) circle (1.7pt) node[below left] {$3$};
  \fill (2,0) circle (1.7pt) node[below right] {$4$};
  \node[courseblue,right] at (2.4,1) {reflection axis};
\end{tikzpicture}
\end{center}
Counterclockwise rotation through $90^\circ$ is the permutation
$x=(1\,2\,3\,4)$, and reflection across the diagonal through vertices
$1$ and $3$ is $r=(2\,4)$. They satisfy
\[
  x^4=1,\qquad r^2=1,\qquad rxr^{-1}=x^{-1}.
\]
The last relation expresses the fact that reflecting a rotation reverses
its direction. Notice that
\[
  x^2=(1\,3)(2\,4),\qquad xr=(1\,2)(3\,4),\qquad
  x^3r=(1\,4)(2\,3),
\]
so the group generated by $x$ and $r$ contains $K$. What is this group?


The subgroup $\langle x,r\rangle\leq S_4$ consists of the four rotations
and four reflections of the square:
\[
\{1,x,x^2,x^3,r,xr,x^2r,x^3r\}.
\]
These eight symmetries are distinct. This is an example of a more
general construction:


\begin{define}[Dihedral group]
	For an integer $n\geq3$, the \emph{dihedral group} $D_{2n}$ is the group of
	symmetries of a regular $n$-gon. It is generated by a rotation $x$ of
	order $n$ and a reflection $y$ of order $2$, with relations
	\[
	x^n=1,\qquad y^2=1,\qquad yxy^{-1}=x^{-1}.
	\]
The last relation is equivalent to $yx=x^{-1}y=x^{n-1}y$.
Our notation records the order of the group in its subscript.
\end{define}

\begin{exercise}
Verify that the relation $yx=x^{-1}y$ lets us write every element as
either $x^i$ or $x^iy$, where $0\leq i<n$. Show that the following
$2n$ elements of $D_{2n}$ are distinct:
\[
  \{1,x,x^2,\ldots,x^{n-1},y,xy,x^2y,\ldots,x^{n-1}y\}.
\]
Geometrically, these are the $n$ rotations and $n$ reflections, so
$|D_{2n}|=2n$.

For a regular hexagon, $n=6$ and $D_{12}$ has six rotations and six
reflections. Label the vertices counterclockwise as shown below.
The blue arrows show two rotations; the dashed lines show two
reflection axes.
\begin{center}
\begin{tikzpicture}[x=1cm,y=1cm,every node/.style={font=\small,align=center}]
  \foreach \offset/\panel in {0/1,3.5/2,7/3,10.5/4} {
    \begin{scope}[xshift=\offset cm]
      \ifnum\panel=1
        \draw[-stealth,courseblue,line width=1pt]
          (0:1.6) arc[start angle=0,end angle=60,radius=1.6];
      \fi
      \ifnum\panel=2
        \draw[-stealth,courseblue,line width=1pt]
          (0:1.6) arc[start angle=0,end angle=120,radius=1.6];
      \fi
      \ifnum\panel=3
        \draw[dashed,ForestGreen,line width=1pt] (180:1.6)--(0:1.6);
      \fi
      \ifnum\panel=4
        \draw[dashed,BrickRed,line width=1pt] (210:1.6)--(30:1.6);
      \fi
      \foreach \vertex in {1,...,6} {
        \coordinate (v\vertex) at ({60*(\vertex-1)}:1.05);
      }
      \draw[thick] (v1)--(v2)--(v3)--(v4)--(v5)--(v6)--cycle;
      \foreach \vertex in {1,...,6} {
        \fill (v\vertex) circle (1.4pt);
        \node[fill=exercisebackground,inner sep=1pt]
          at ({60*(\vertex-1)}:1.33) {$\vertex$};
      }
      \ifnum\panel=1
        \node[courseblue] at (0,-1.75) {$x$: $60^\circ$ rotation};
      \fi
      \ifnum\panel=2
        \node[courseblue] at (0,-1.75) {$x^2$: $120^\circ$ rotation};
      \fi
      \ifnum\panel=3
        \node[ForestGreen] at (0,-1.75) {$y$: reflection\\vertex axis};
      \fi
      \ifnum\panel=4
        \node[BrickRed] at (0,-1.75) {$xy$: reflection\\edge-midpoint axis};
      \fi
    \end{scope}
  }
\end{tikzpicture}
\end{center}
For these labels, the symmetries are
\[
\begin{aligned}
  x&=(1\,2\,3\,4\,5\,6), & x^2&=(1\,3\,5)(2\,4\,6),\\
  y&=(2\,6)(3\,5), & xy&=(1\,2)(3\,6)(4\,5).
\end{aligned}
\]
The axis of $xy$ passes through the midpoints of the edges
$\{1,2\}$ and $\{4,5\}$.
\end{exercise}

The following pictures illustrate why the square symmetry group
contains $K$.

\begin{remark}[Why the square symmetry group contains $K$]
Return to the square, where $x=(1\,2\,3\,4)$ and $r=(2\,4)$.
Its symmetry group is $\langle x,r\rangle\cong D_8$. Direct computation
gives
\[
  x^2=(1\,3)(2\,4),\qquad
  xr=(1\,2)(3\,4),\qquad
  x^3r=(1\,4)(2\,3).
\]
Together with the identity, these give all four elements of the kernel:
\[
  K=\{1,(1\,2)(3\,4),(1\,3)(2\,4),(1\,4)(2\,3)\}
   =\{1,x^2,xr,x^3r\}\subseteq\langle x,r\rangle.
\]
Geometrically, the three nonidentity elements are the half-turn and
the two reflections through opposite edge midpoints, as pictured below.
\begin{center}
\begin{tikzpicture}[x=1cm,y=1cm,every node/.style={font=\small,align=center}]
  \foreach \offset/\panel in {0/1,4.5/2,9/3} {
    \begin{scope}[xshift=\offset cm]
      \ifnum\panel=1
        \draw[-stealth,courseblue,line width=1pt]
          (45:1.65) arc[start angle=45,end angle=225,radius=1.65];
      \fi
      \ifnum\panel=2
        \draw[dashed,ForestGreen,line width=1pt] (0,-1.4)--(0,1.4);
      \fi
      \ifnum\panel=3
        \draw[dashed,BrickRed,line width=1pt] (-1.4,0)--(1.4,0);
      \fi
      \draw[thick] (-0.8,-0.8)--(-0.8,0.8)--(0.8,0.8)--(0.8,-0.8)--cycle;
      \fill (0.8,0.8) circle (1.5pt) node[above right] {$1$};
      \fill (-0.8,0.8) circle (1.5pt) node[above left] {$2$};
      \fill (-0.8,-0.8) circle (1.5pt) node[below left] {$3$};
      \fill (0.8,-0.8) circle (1.5pt) node[below right] {$4$};
      \ifnum\panel=1
        \node[courseblue] at (0,-1.85)
          {$x^2=(1\,3)(2\,4)$\\$180^\circ$ half-turn};
      \fi
      \ifnum\panel=2
        \node[ForestGreen] at (0,-1.85)
          {$xr=(1\,2)(3\,4)$\\vertical-axis reflection};
      \fi
      \ifnum\panel=3
        \node[BrickRed] at (0,-1.85)
          {$x^3r=(1\,4)(2\,3)$\\horizontal-axis reflection};
      \fi
    \end{scope}
  }
\end{tikzpicture}
\end{center}
\end{remark}

Moreover, $\varphi(x)=\varphi(r)=(1\,3)$, so
$\langle x,r\rangle\leq\varphi^{-1}(\langle(1\,3)\rangle)$.
The preimage has order $|K|\cdot2=8$ by the correspondence theorem,
and hence $\langle x,r\rangle=\varphi^{-1}(\langle(1\,3)\rangle)=H_1$.

For $H_2$ and $H_3$, relabel the vertices as
shown below. The opposite vertex pairs now form $\Pi_1$ and $\Pi_3$,
respectively. The curved arrows indicate counterclockwise quarter-turns;
the green dashed lines are diagonal reflection axes.
\begin{center}
\begin{tikzpicture}[x=1cm,y=1cm,every node/.style={font=\small,align=center}]
  \begin{scope}
    \node[courseblue,font=\bfseries] at (0,2.4) {$H_2$};
    \draw[-stealth,courseblue,line width=1pt]
      (45:1.65) arc[start angle=45,end angle=135,radius=1.65];
    \node[courseblue] at (0,1.9) {$x_2$};
    \draw[dashed,ForestGreen,line width=1pt] (-1.3,-1.3)--(1.3,1.3);
    \node[ForestGreen,right] at (1.35,1.3) {$r_2$};
    \draw[thick] (-0.9,-0.9)--(-0.9,0.9)--(0.9,0.9)--(0.9,-0.9)--cycle;
    \fill (0.9,0.9) circle (1.5pt) node[above] {$1$};
    \fill (-0.9,0.9) circle (1.5pt) node[above] {$3$};
    \fill (-0.9,-0.9) circle (1.5pt) node[below] {$2$};
    \fill (0.9,-0.9) circle (1.5pt) node[below] {$4$};
    \node[courseblue] at (0,-1.45) {opposite pairs: $\Pi_1$};
    \node at (0,-2) {$x_2=(1\,3\,2\,4),\quad r_2=(3\,4)$};
  \end{scope}
  \begin{scope}[xshift=7cm]
    \node[BrickRed,font=\bfseries] at (0,2.4) {$H_3$};
    \draw[-stealth,BrickRed,line width=1pt]
      (45:1.65) arc[start angle=45,end angle=135,radius=1.65];
    \node[BrickRed] at (0,1.9) {$x_3$};
    \draw[dashed,ForestGreen,line width=1pt] (-1.3,-1.3)--(1.3,1.3);
    \node[ForestGreen,right] at (1.35,1.3) {$r_3$};
    \draw[thick] (-0.9,-0.9)--(-0.9,0.9)--(0.9,0.9)--(0.9,-0.9)--cycle;
    \fill (0.9,0.9) circle (1.5pt) node[above] {$1$};
    \fill (-0.9,0.9) circle (1.5pt) node[above] {$2$};
    \fill (-0.9,-0.9) circle (1.5pt) node[below] {$4$};
    \fill (0.9,-0.9) circle (1.5pt) node[below] {$3$};
    \node[BrickRed] at (0,-1.45) {opposite pairs: $\Pi_3$};
    \node at (0,-2) {$x_3=(1\,2\,4\,3),\quad r_3=(2\,3)$};
  \end{scope}
\end{tikzpicture}
\end{center}
The induced permutations of the three pairings are
\[
  \varphi(x_2)=\varphi(r_2)=(2\,3)=ab,\qquad
  \varphi(x_3)=\varphi(r_3)=(1\,2)=a^2b.
\]
Each square symmetry group has order $8$, as does the corresponding
preimage. Therefore
\[
\begin{aligned}
  H_2&=\langle x_2,r_2\rangle=\varphi^{-1}(\langle ab\rangle),\\
  H_3&=\langle x_3,r_3\rangle=\varphi^{-1}(\langle a^2b\rangle).
\end{aligned}
\]
Both are copies of $D_8$ inside $S_4$ containing $K$.

\begin{exercise}
Use the correspondence theorem to explain why none of the three
subgroups $H_1,H_2,H_3$ is normal in $S_4$.
\end{exercise}

\end{document}
